% ============================================================================
% text_formatting.tex - Text Formatting in LaTeX
% ============================================================================
% Demonstrates font families, styles, sizes, colors, lists, verbatim,
% special characters, text alignment, and footnotes.
%
% Compilation:
%   pdflatex text_formatting.tex
% ============================================================================

\documentclass[12pt,a4paper]{article}

% ============================================================================
% PACKAGE IMPORTS
% ============================================================================

\usepackage[utf8]{inputenc}      % UTF-8 input encoding
\usepackage[T1]{fontenc}         % Font encoding for European characters
\usepackage{lmodern}             % Latin Modern font (scalable CM replacement)

\usepackage[margin=1in]{geometry}
\usepackage[english]{babel}

\usepackage{xcolor}              % Extended color support
\usepackage{listings}            % Source code formatting
\usepackage{verbatim}            % Verbatim text environments

\usepackage{hyperref}
\hypersetup{
    colorlinks=true,
    linkcolor=blue,
    urlcolor=cyan,
    pdftitle={Text Formatting in LaTeX},
    pdfauthor={LaTeX Student}
}

% ============================================================================
% CUSTOM COLORS
% ============================================================================

\definecolor{myblue}{RGB}{0, 70, 140}        % Custom deep blue
\definecolor{mygold}{RGB}{200, 150, 0}       % Custom gold
\definecolor{mygreen}{RGB}{20, 120, 50}      % Custom forest green
\definecolor{highlight}{RGB}{255, 245, 150}  % Pale yellow for highlighting

% ============================================================================
% LISTINGS CONFIGURATION
% ============================================================================

\lstset{
    basicstyle=\ttfamily\small,
    keywordstyle=\color{myblue}\bfseries,
    commentstyle=\color{mygreen}\itshape,
    stringstyle=\color{red!70!black},
    numbers=left,
    numberstyle=\tiny\color{gray},
    numbersep=8pt,
    frame=single,
    framecolor=gray!40,
    rulecolor=\color{gray!40},
    breaklines=true,
    showstringspaces=false,
    tabsize=4
}

% ============================================================================
% DOCUMENT METADATA
% ============================================================================

\title{\textbf{Text Formatting in \LaTeX}\\[0.5em]
       \large A Practical Reference}
\author{LaTeX Student}
\date{\today}

% ============================================================================
% DOCUMENT BODY
% ============================================================================

\begin{document}

\maketitle
\tableofcontents
\newpage

% ============================================================================
% SECTION 1: FONT FAMILIES
% ============================================================================

\section{Font Families}
\label{sec:families}

\LaTeX{} provides three fundamental font families, each suited to different
contexts in a document.

\subsection{Serif (Roman) — \texttt{\textbackslash textrm}}

The default family. Serifs are the small strokes at the ends of letterforms,
which aid readability in body text.

\textrm{This text uses the Roman (serif) font family. Serifs guide the eye
along horizontal lines, making long passages comfortable to read.}

\subsection{Sans Serif — \texttt{\textbackslash textsf}}

Sans-serif text omits serifs, giving a clean, modern appearance well-suited
to headings, captions, and on-screen display.

\textsf{This text uses the sans-serif font family. It looks clean and
contemporary, often used in slide presentations and user interfaces.}

\subsection{Monospace (Typewriter) — \texttt{\textbackslash texttt}}

Every character occupies the same horizontal width, making monospace fonts
ideal for source code, file paths, and terminal output.

\texttt{This text uses the monospace (typewriter) font family. Each character
takes equal horizontal space: i m w l all line up precisely.}

% ============================================================================
% SECTION 2: FONT STYLES AND EMPHASIS
% ============================================================================

\section{Font Styles and Emphasis}
\label{sec:styles}

\subsection{Basic Styles}

\begin{center}
\begin{tabular}{lll}
\hline
\textbf{Command} & \textbf{Example} & \textbf{Use case} \\
\hline
\verb|\textbf{...}| & \textbf{Bold text} & Key terms, strong emphasis \\
\verb|\textit{...}| & \textit{Italic text} & Titles, foreign words \\
\verb|\textsl{...}| & \textsl{Slanted text} & Lighter emphasis than italic \\
\verb|\textsc{...}| & \textsc{Small Capitals} & Acronyms, proper names \\
\verb|\underline{...}| & \underline{Underlined} & Use sparingly \\
\hline
\end{tabular}
\end{center}

\subsection{Combining Styles}

Styles can be nested, though excessive combinations harm readability:

\medskip
\textbf{Bold} combined with \textbf{\textit{italic}} produces bold italic.

\texttt{\textbf{Bold monospace}} is useful for highlighting code keywords.

\textit{\textsc{Italic small caps}} works for special terminology.

\medskip
\noindent
The semantic commands \verb|\emph{...}| toggle between roman and italic
depending on context: in normal text \emph{this is italic}, but inside
italic text \textit{this word \emph{reverts} to roman}.

% ============================================================================
% SECTION 3: FONT SIZES
% ============================================================================

\section{Font Sizes}
\label{sec:sizes}

\LaTeX{} provides ten relative size commands, scaled relative to the
base font size declared in \verb|\documentclass|:

\medskip
{\tiny        \texttt{tiny}        — 5pt at 10pt base}

{\scriptsize  \texttt{scriptsize}  — 7pt at 10pt base}

{\footnotesize\texttt{footnotesize}— 8pt at 10pt base}

{\small       \texttt{small}       — 9pt at 10pt base}

{\normalsize  \texttt{normalsize}  — 10pt at 10pt base (default)}

{\large       \texttt{large}       — 12pt at 10pt base}

{\Large       \texttt{Large}       — 14pt at 10pt base}

{\LARGE       \texttt{LARGE}       — 17pt at 10pt base}

{\huge        \texttt{huge}        — 20pt at 10pt base}

{\Huge        \texttt{Huge}        — 25pt at 10pt base}

\medskip
\noindent
Note: braces \verb|{ }| limit the scope of a size declaration, preventing
it from affecting the rest of the document.

% ============================================================================
% SECTION 4: COLORS WITH xcolor
% ============================================================================

\section{Colors with \texttt{xcolor}}
\label{sec:colors}

\subsection{Predefined Colors}

The \texttt{xcolor} package provides 68 predefined color names from the
\texttt{dvipsnames} option (loaded by default in many setups):

\medskip
\textcolor{red}{Red text} \quad
\textcolor{blue}{Blue text} \quad
\textcolor{green}{Green text} \quad
\textcolor{cyan}{Cyan text} \quad
\textcolor{magenta}{Magenta text} \quad
\textcolor{orange}{Orange text} \quad
\textcolor{violet}{Violet text} \quad
\textcolor{brown}{Brown text}

\subsection{Custom Colors with \texttt{\textbackslash definecolor}}

\begin{verbatim}
\definecolor{myblue}{RGB}{0, 70, 140}     % RGB model (0-255 each)
\definecolor{mygold}{RGB}{200, 150, 0}
\definecolor{mygreen}{RGB}{20, 120, 50}
\end{verbatim}

Result:
\textcolor{myblue}{This is custom deep blue.} \quad
\textcolor{mygold}{This is custom gold.} \quad
\textcolor{mygreen}{This is custom forest green.}

\medskip
Colors can also be defined with other models:

\definecolor{warm}{HTML}{CC4400}      % Hex notation
\definecolor{pale}{gray}{0.75}        % Grayscale (0=black, 1=white)

\textcolor{warm}{HTML color \#CC4400} \quad
\textcolor{pale}{Gray (0.75)}

\subsection{Colored Backgrounds with \texttt{\textbackslash colorbox}}

\colorbox{highlight}{Highlighted text on a pale yellow background.}

\medskip
\colorbox{myblue}{\textcolor{white}{\textbf{White text on dark blue.}}}

\medskip
For boxes with a colored frame \emph{and} background, use
\verb|\fcolorbox{frame}{bg}{text}|:

\fcolorbox{myblue}{highlight}{Framed and highlighted text.}

\subsection{Color Mixing}

\texttt{xcolor} supports inline color expressions:

\textcolor{red!50!blue}{50\% red + 50\% blue = purple}

\textcolor{green!70!black}{70\% green + 30\% black = dark green}

\textcolor{blue!30}{30\% blue + 70\% white = light blue}

% ============================================================================
% SECTION 5: LIST ENVIRONMENTS
% ============================================================================

\section{List Environments}
\label{sec:lists}

\subsection{Unordered Lists (\texttt{itemize})}

Use \texttt{itemize} for collections where order does not matter:

\begin{itemize}
    \item The \LaTeX{} kernel provides core functionality
    \item Packages extend \LaTeX{} for specialized tasks
    \item Document classes define the overall layout
    \item The \texttt{CTAN} archive hosts thousands of packages
\end{itemize}

\subsection{Ordered Lists (\texttt{enumerate})}

Use \texttt{enumerate} for sequences, steps, or ranked items:

\begin{enumerate}
    \item Install a \TeX{} distribution (TeX Live or MiKTeX)
    \item Choose an editor (TeXstudio, VS Code, Overleaf)
    \item Create a \texttt{.tex} source file
    \item Run \texttt{pdflatex} to compile
    \item View the resulting PDF
\end{enumerate}

\subsection{Description Lists (\texttt{description})}

Use \texttt{description} for glossaries, symbol tables, or labeled items:

\begin{description}
    \item[\texttt{pdflatex}] The most common \LaTeX{} engine; produces PDF directly
    \item[\texttt{xelatex}] Supports Unicode input and OpenType fonts natively
    \item[\texttt{lualatex}] Embeds Lua scripting in \LaTeX; highly extensible
    \item[\texttt{biber}]    Bibliography processor for BibLaTeX (replaces BibTeX)
    \item[\texttt{latexmk}]  Automates multi-pass compilation for cross-references
\end{description}

\subsection{Nested Lists}

Lists can nest up to four levels deep (default labels change automatically):

\begin{enumerate}
    \item Programming paradigms
        \begin{itemize}
            \item Imperative
                \begin{enumerate}
                    \item Procedural (C, Pascal)
                    \item Object-oriented (Java, Python)
                \end{enumerate}
            \item Declarative
                \begin{itemize}
                    \item Functional (Haskell, Lisp)
                    \item Logic (Prolog)
                \end{itemize}
        \end{itemize}
    \item Type systems
        \begin{itemize}
            \item Static typing (C, Rust, Haskell)
            \item Dynamic typing (Python, JavaScript, Ruby)
        \end{itemize}
\end{enumerate}

% ============================================================================
% SECTION 6: VERBATIM AND CODE LISTINGS
% ============================================================================

\section{Verbatim and Code Listings}
\label{sec:verbatim}

\subsection{Inline Verbatim}

Use \verb|\verb| for short inline code: the command \verb|\textbf{}| makes
text bold. File paths such as \verb|/usr/local/bin/pdflatex| stay literal.

\subsection{The \texttt{verbatim} Environment}

The \texttt{verbatim} environment preserves whitespace and suppresses
all \LaTeX{} interpretation:

\begin{verbatim}
$ pdflatex --interaction=nonstopmode document.tex
$ biber document
$ pdflatex document.tex
$ pdflatex document.tex
\end{verbatim}

\subsection{Syntax-Highlighted Listings}

The \texttt{listings} package adds syntax highlighting with line numbers:

\begin{lstlisting}[language=Python, caption={Sieve of Eratosthenes}]
def sieve(limit):
    """Return all prime numbers up to limit."""
    is_prime = [True] * (limit + 1)
    is_prime[0] = is_prime[1] = False

    for i in range(2, int(limit**0.5) + 1):
        if is_prime[i]:
            # Mark all multiples as composite
            for j in range(i * i, limit + 1, i):
                is_prime[j] = False

    return [n for n, prime in enumerate(is_prime) if prime]

primes = sieve(50)
print(f"Primes up to 50: {primes}")
\end{lstlisting}

% ============================================================================
% SECTION 7: SPECIAL CHARACTERS AND DASHES
% ============================================================================

\section{Special Characters and Dashes}
\label{sec:special}

\subsection{Reserved Characters}

These ten characters have special meaning in \LaTeX{} and must be escaped:

\medskip
\begin{center}
\begin{tabular}{clc}
\hline
\textbf{Character} & \textbf{Command} & \textbf{Usage} \\
\hline
\#  & \verb|\#|  & Parameter marker in macros \\
\$  & \verb|\$|  & Math mode delimiter \\
\%  & \verb|\%|  & Comment character \\
\&  & \verb|\&|  & Column separator in tables \\
\_  & \verb|\_|  & Subscript in math mode \\
\{  & \verb|\{|  & Begin group \\
\}  & \verb|\}|  & End group \\
\~{} & \verb|\~{}| & Non-breaking space (tilde) \\
\^{} & \verb|\^{}| & Superscript in math mode \\
\textbackslash & \verb|\textbackslash| & Macro prefix \\
\hline
\end{tabular}
\end{center}

\subsection{Dash Types}

\LaTeX{} uses different numbers of hyphens for three distinct purposes:

\begin{description}
    \item[Hyphen \texttt{-}] Used in compound words: daughter-in-law, twenty-one
    \item[En dash \texttt{--}] Used for ranges and connections: pages 10--25, London--Paris route
    \item[Em dash \texttt{---}] Used for parenthetical remarks---like this---in prose
\end{description}

\subsection{Quotation Marks}

\LaTeX{} uses backticks for opening quotes, straight quotes for closing:

\begin{description}
    \item[Single quotes] \verb|`text'| produces `text'
    \item[Double quotes] \verb|``text''| produces ``text''
    \item[Nested] \verb|``She said `hello' to me.''| produces ``She said `hello' to me.''
\end{description}

\subsection{Other Typographic Symbols}

\copyright{} 2024 LaTeX Foundation \quad
\textregistered{} Registered Trademark \quad
\texttrademark{} Trademark \quad
\pounds 42 \quad
\S\,3.1 (section sign) \quad
\P{} (pilcrow) \quad
\dag{} (dagger) \quad
\ddag{} (double dagger)

\medskip
\noindent Ellipsis: use \verb|\ldots| to get proper spacing\ldots{} not three dots...

% ============================================================================
% SECTION 8: TEXT ALIGNMENT
% ============================================================================

\section{Text Alignment}
\label{sec:alignment}

By default, \LaTeX{} fully justifies text, stretching spaces to align both
left and right margins. Three environments override this:

\subsection{Centered Text}

\begin{center}
    \textbf{Centered Heading}\\
    This paragraph is centered on the page.\\
    Each line is independently centered regardless of its length.
\end{center}

\subsection{Left-Aligned (Ragged Right)}

\begin{flushleft}
    This paragraph is left-aligned with a ragged right margin.
    No interword spacing adjustments are made to justify lines.
    This style is common in letters and informal documents.
\end{flushleft}

\subsection{Right-Aligned (Ragged Left)}

\begin{flushright}
    This paragraph is right-aligned with a ragged left margin.\\
    Often used for dates, addresses, and captions in some styles.\\
    \textit{—Anonymous}
\end{flushright}

\subsection{The \texttt{\textbackslash hfill} and \texttt{\textbackslash vfill} Commands}

\texttt{\textbackslash hfill} inserts a flexible horizontal space that
fills all available width:

\noindent Author Name \hfill \today

\medskip
\noindent Left \hfill Center \hfill Right

% ============================================================================
% SECTION 9: FOOTNOTES
% ============================================================================

\section{Footnotes}
\label{sec:footnotes}

Footnotes\footnote{A footnote appears at the bottom of the page, numbered
automatically. Keep them concise and non-essential to the main argument.}
are placed with \verb|\footnote{text}| immediately after the word or phrase
they annotate.

Multiple footnotes on a page are numbered sequentially.\footnote{This is the
second footnote on this page. Notice the automatic numbering.} The numbering
resets at the start of each new chapter in \texttt{book} and \texttt{report}
class documents.

For tables and other environments where \verb|\footnote| cannot be used
directly, use the two-step approach: place \verb|\footnotemark| in the table
cell, then \verb|\footnotetext{...}| after the table.\footnotemark

\footnotetext{This footnote was placed using \texttt{\textbackslash footnotemark}
inside the paragraph and \texttt{\textbackslash footnotetext} outside, because
some environments do not support \texttt{\textbackslash footnote} directly.}

Custom footnote numbers are possible with \verb|\footnote[n]{text}|: see
footnote~\ref{fn:custom}.\footnote[99]{This is footnote number 99, set
with \texttt{\textbackslash footnote[99]\{...\}}.\label{fn:custom}}

% ============================================================================
% SUMMARY
% ============================================================================

\section*{Summary}

This document has demonstrated:
\begin{itemize}
    \item The three font families: roman (\verb|\textrm|), sans-serif (\verb|\textsf|),
          and monospace (\verb|\texttt|)
    \item Font styles: bold, italic, slanted, small caps, and combinations
    \item Ten relative font sizes from \verb|\tiny| to \verb|\Huge|
    \item Color commands from the \texttt{xcolor} package: predefined colors,
          custom \verb|\definecolor|, \verb|\colorbox|, and color mixing
    \item Three list environments: \texttt{itemize}, \texttt{enumerate},
          \texttt{description}, and nested combinations
    \item Verbatim text with \verb|\verb| and the \texttt{verbatim} environment,
          plus syntax-highlighted listings with the \texttt{listings} package
    \item Special characters, dash types, and quotation marks
    \item Alignment environments: \texttt{center}, \texttt{flushleft},
          \texttt{flushright}
    \item Footnotes with automatic and manual numbering
\end{itemize}

\end{document}

% ============================================================================
% COMPILATION NOTES
% ============================================================================
% Single pass is sufficient for this document:
%   pdflatex text_formatting.tex
%
% For documents that also include cross-references or bibliography:
%   latexmk -pdf text_formatting.tex
% ============================================================================
